% ============================================================================
% Companion appendix for gw_section_body.tex.
% When merging into Main.tex, place this AFTER the host's \appendix switch,
% alongside the existing appendices. It is the compressed, sign-checked
% replacement for the level-set subsection of the original draft, retargeted
% at simulators. Flat space, natural units, host conventions.
% ============================================================================

\section{Implementing finite-rate conversion in merger simulations}
\label{sec:finite_rate_implementation}

The finite-rate prescription of Sec.~\ref{sec:gw} can be added to a hydrodynamic merger simulation with one additional advected field.
Let $\Phi(\boldsymbol{x},t)$ be a marker field whose zero surface is the phase boundary, with $\Phi>0$ in quark matter and $\Phi<0$ in nuclear matter, so that the conversion indicator is $Y_q=\Theta(\Phi)$.
Only the zero surface and the sign of $\Phi$ are physical; its magnitude may be chosen for numerical convenience, e.g.\ by reinitializing $\Phi$ to the signed distance from the boundary.

The boundary moves with the local fluid velocity $\boldsymbol{v}$ plus the conversion speed $u(0^-)$ along the normal $\widehat{\boldsymbol{n}}=-\boldsymbol{\nabla}\Phi/|\boldsymbol{\nabla}\Phi|$, which points from the quark phase into the nuclear phase.
Requiring the zero surface to follow this velocity gives
\begin{equation}
\partial_t\Phi
+\boldsymbol{v}\boldsymbol{\cdot}\boldsymbol{\nabla}\Phi
=u(0^-)\,\bigl|\boldsymbol{\nabla}\Phi\bigr|\,,
\label{eq:marker_transport}
\end{equation}
a level-set equation whose source term converts matter at the calculated speed.
Here $\boldsymbol{v}$ is the fluid velocity on the nuclear side, because $u(0^-)$ is defined relative to the nuclear matter flowing into the front; the boundary speed relative to the quark side differs by the density ratio through baryon-flux conservation (Sec.~\ref{sec:speed}).
The speed $u(0^-)$ is evaluated pointwise on the boundary from the local values of $n_B(0^-)-n_B^{\rm crit}(T)$ and $T$, using the analytic boundary-speed formula of Sec.~\ref{sec:speed} or a tabulation of the numerical solutions.

Two limitations of the prescription mirror the physics.
First, $Y_q=\Theta(\Phi)$ builds in hysteresis: where the front has passed, a later pressure drop below $p_\text{crit}$ does not undo the conversion unless a reverse-conversion speed is supplied on the quark side of the boundary.
Second, Eq.~\eqref{eq:marker_transport} transports existing boundaries only; where $P>p_\text{crit}$ but no boundary is present, conversion requires nucleation, which must be seeded explicitly.
